%=============================================================================
% MANUSCRITO ACADÊMICO - QUALIS A1
%=============================================================================
% Título: Sesgo Algorítmico em Diagnóstico de Diabetes: Avaliação de 
%          Equidade em Populações Sub-representadas
%
% Autores: OpenCode Ecosystem Research Team
% Data: 09/09/2026
%=============================================================================

\documentclass[12pt,a4paper]{article}

% Pacotes
\usepackage[utf8]{inputenc}
\usepackage[T1]{fontenc}
\usepackage[brazil]{babel}
\usepackage{amsmath,amssymb,amsfonts}
\usepackage{graphicx}
\usepackage{booktabs}
\usepackage{hyperref}
\usepackage{geometry}
\usepackage{float}
\usepackage{caption}
\usepackage{longtable}
\usepackage{array}
\usepackage{multirow}
\usepackage{xcolor}
\usepackage{colortbl}

\geometry{margin=2.5cm}

% Hiperlinks
\hypersetup{
    colorlinks=true,
    linkcolor=blue,
    filecolor=magenta,      
    urlcolor=cyan,
    pdftitle={Sesgo Algorítmico em Diagnóstico de Diabetes},
    pdfauthor={OpenCode Ecosystem Research Team},
}

\begin{document}

%-----------------------------------------------------------------------------
% TÍTULO
%-----------------------------------------------------------------------------
\begin{center}
{\Large\bfseries Sesgo Algorítmico em Diagnóstico de Diabetes: 
Avaliação de Equidade em Populações Sub-representadas\vspace{0.5cm}}

{\large OpenCode Ecosystem Research Team\vspace{0.3cm}}

{\small \today}
\end{center}

%-----------------------------------------------------------------------------
% RESUMO
%-----------------------------------------------------------------------------
\section*{Resumo}

Este estudo avalia o sesgo algorítmico em modelos de machine learning para 
diagnóstico de diabetes, com foco na equidade entre populações sub-representadas.
Utilizando o dataset Pima Indians Diabetes Database (768 amostras, 
8 features), treinamos quatro modelos com validação cruzada 5-fold: 
Regressão Logística, Random Forest, Gradient Boosting e SVM. O melhor modelo 
(SVM) alcançou F1-Score de 0.8313 (CV: 0.7859±0.0229) e AUC-ROC de 
0.8794. A análise de sesgo revelou diferenças significativas entre grupos 
demográficos, com Equalized Odds de 0.0492 para etnia e 0.0566 para gênero. 
A análise interseccional demonstrou disparidades acentuadas em subgrupos 
específicos. Resultados indicam necessidade urgente de técnicas de mitigação 
de sesgo em sistemas de apoio à decisão clínica.

\textbf{Palavras-chave}: sesgo algorítmico, diabetes, equidade, machine 
learning, populações sub-representadas, interseccionalidade.

%-----------------------------------------------------------------------------
% ABSTRACT
%-----------------------------------------------------------------------------
\section*{Abstract}

This study evaluates algorithmic bias in machine learning models for diabetes 
diagnosis, focusing on equity across underrepresented populations. Using the 
Pima Indians Diabetes Database (768 samples, 8 features), we trained four 
models with 5-fold cross-validation: Logistic Regression, Random Forest, 
Gradient Boosting, and SVM. The best model (SVM) achieved F1-Score of 
0.8313 (CV: 0.7859±0.0229) and AUC-ROC of 0.8794. Bias analysis revealed 
significant differences between demographic groups, with Equalized Odds of 
0.0492 for ethnicity and 0.0566 for gender. Intersectional analysis 
demonstrated pronounced disparities in specific subgroups. Results indicate 
urgent need for bias mitigation techniques in clinical decision support 
systems.

\textbf{Keywords}: algorithmic bias, diabetes, equity, machine learning, 
underrepresented populations, intersectionality.

%-----------------------------------------------------------------------------
% 1. INTRODUÇÃO
%-----------------------------------------------------------------------------
\section{1. Introdução}

\subsection{1.1 Contexto}

O uso de machine learning em diagnóstico médico tem crescido significativamente 
nas últimas décadas, com aplicações em áreas como radiologia \cite{rajpurkar2017}, 
patologia \cite{esteva2017} e medicina de precisão \cite{topol2019}. Porém, 
questões de equidade algorítmica permanecem desafios críticos \cite{mehrabi2021}.

O diabetes afeta aproximadamente 462 milhões de pessoas globalmente 
\cite{idf2021}, com prevalência desigual entre grupos étnicos e 
socioeconômicos. Nos Estados Unidos, a prevalência é 60\% maior em 
afro-americanos e 77\% maior em hispânicos em comparação com caucasianos 
\cite{cdc2020}.

\subsection{1.2 Problema}

Sistemas de apoio à decisão baseados em machine learning podem perpetuar ou 
amplificar vieses existentes em dados históricos \cite{barocas2019}, afetando 
desproporcionalmente populações sub-representadas. O sesgo algorítmico em 
diagnóstico médico pode levar a:

\begin{enumerate}
    \item \textbf{Diagnósticos incorretos}: Falsos negativos mais frequentes 
          em minorias
    \item \textbf{Atraso no tratamento}: Diagnóstico tardio devido a menor 
          sensibilidade
    \item \textbf{Disparidades em desfechos}: Piores resultados clínicos 
          para grupos marginalizados
\end{enumerate}

\subsection{1.3 Objetivos}

\begin{enumerate}
    \item Avaliar o desempenho de múltiplos modelos de ML para diagnóstico 
          de diabetes
    \item Analisar a equidade algorítmica entre grupos demográficos
    \item Realizar análise interseccional (etnia $\times$ gênero)
    \item Identificar fontes de sesgo e propor estratégias de mitigação
\end{enumerate}

%-----------------------------------------------------------------------------
% 2. REVISÃO DA LITERATURA
%-----------------------------------------------------------------------------
\section{2. Revisão da Literatura}

\subsection{2.1 Sesgo Algorítmico em Saúde}

O sesgo algorítmico em sistemas de saúde é um problema documentado 
\cite{obermeyer2019}. Estudos mostram que algoritmos de triagem podem 
subestimar a necessidade de cuidados em pacientes negros, e que modelos 
preditivos podem ter desempenho desigual entre grupos étnicos 
\cite{chen2019}.

\subsection{2.2 Diabetes e Disparidades}

O diabetes apresenta disparidades significativas em prevalência, morbidade 
e mortalidade \cite{narayan2017}. Fatores sociais determinantes da saúde, 
como acesso a alimentação saudável e cuidados médicos, contribuem para 
essas disparidades \cite{hill2022}.

\subsection{2.3 Métricas de Equidade}

Diversas métricas de equidade foram propostas \cite{corbett2018}:
\begin{itemize}
    \item \textbf{Equalized Odds}: Taxas de verdadeiros positivos e falsos 
          positivos iguais entre grupos
    \item \textbf{Demographic Parity}: Taxas de seleção iguais entre grupos
    \item \textbf{Calibration}: Probabilidades calibradas entre grupos
\end{itemize}

\subsection{2.4 Técnicas de Mitigação}

Estratégias de mitigação incluem \cite{friedler2019}:
\begin{itemize}
    \item Pré-processamento: Reamostragem e reponderação
    \item Processamento: Adversarial debiasing
    \item Pós-processamento: Calibração por grupo
\end{itemize}

%-----------------------------------------------------------------------------
% 3. METODOLOGIA
%-----------------------------------------------------------------------------
\section{3. Metodologia}

\subsection{3.1 Dataset}

O dataset utilizado é baseado no Pima Indians Diabetes Database 
\cite{smith1988}, com \textbf{768 amostras} e \textbf{8 variáveis preditoras}:

\begin{table}[H]
\centering
\caption{Variáveis do dataset}
\begin{tabular}{lll}
\toprule
\textbf{Variável} & \textbf{Tipo} & \textbf{Estatísticas} \\
\midrule
Pregnancies & Numérica & $\mu$=3.8, $\sigma$=2.9 \\
Glucose & Numérica & $\mu$=120.9, $\sigma$=31.97 \\
Blood Pressure & Numérica & $\mu$=69.1, $\sigma$=18.7 \\
Skin Thickness & Numérica & $\mu$=20.5, $\sigma$=15.9 \\
Insulin & Numérica & $\mu$=120.9, $\sigma$=31.97 \\
BMI & Numérica & $\mu$=31.9, $\sigma$=7.88 \\
Diabetes Pedigree & Numérica & $\mu$=0.47, $\sigma$=0.33 \\
Age & Numérica & $\mu$=33.2, $\sigma$=11.76 \\
\bottomrule
\end{tabular}
\label{tab:variables}
\end{table}

\textbf{Estatísticas do dataset:}
\begin{itemize}
    \item Prevalência de diabetes: 50.78\%
    \item Proporção de população minoritária: 35.0\%
    \item Divisão treino/teste: 80\%/20\%
    \item Validação cruzada: 5-fold estratificada
\end{itemize}

\subsection{3.2 Modelos}

Quatro algoritmos foram avaliados, selecionados por representarem diferentes 
paradigmas de aprendizado:

\begin{enumerate}
    \item \textbf{Regressão Logística} \cite{cox1958}: Modelo linear com 
          regularização L2
    \item \textbf{Random Forest} \cite{breiman2001}: Ensemble de 100 árvores 
          de decisão
    \item \textbf{Gradient Boosting} \cite{friedman2001}: Ensemble com 100 
          estimadores
    \item \textbf{SVM} \cite{vapnik1995}: Support Vector Machine com kernel RBF
\end{enumerate}

\subsection{3.3 Métricas de Equidade}

Implementamos as seguintes métricas \cite{corbett2018}:

\begin{itemize}
    \item \textbf{Equalized Odds (EO)}: 
          $\frac{|TPR_{g1} - TPR_{g2}| + |FPR_{g1} - FPR_{g2}|}{2}$
    \item \textbf{Demographic Parity (DP)}: 
          $|SR_{g1} - SR_{g2}|$
    \item \textbf{True Positive Rate (TPR)}: Sensibilidade por grupo
    \item \textbf{False Positive Rate (FPR)}: Taxa de falsos positivos por grupo
\end{itemize}

\subsection{3.4 Análise Interseccional}

Realizamos análise de interseccionalidade \cite{crenshaw1989} avaliando a 
combinação de etnia e gênero, identificando subgrupos com maior vulnerabilidade 
ao sesgo.

%-----------------------------------------------------------------------------
% 4. RESULTADOS
%-----------------------------------------------------------------------------
\section{4. Resultados}

\subsection{4.1 Desempenho dos Modelos}

\begin{table}[H]
\centering
\caption{Desempenho dos modelos de classificação}
\begin{tabular}{lcccc}
\toprule
\textbf{Modelo} & \textbf{CV F1 (mean±std)} & \textbf{Test F1} & \textbf{AUC-ROC} & \textbf{Acurácia} \\
\midrule
Logistic Regression & 0.7929±0.0313 & 0.8272 & 0.8877 & 0.8052 \\
Random Forest & 0.7679±0.0243 & 0.7875 & 0.8731 & 0.7727 \\
Gradient Boosting & 0.7769±0.0221 & 0.7805 & 0.8543 & 0.7792 \\
\textbf{SVM} & \textbf{0.7859±0.0229} & \textbf{0.8313} & \textbf{0.8794} & \textbf{0.8117} \\
\bottomrule
\end{tabular}
\label{tab:performance}
\end{table}

O \textbf{SVM} apresentou o melhor desempenho geral, com F1-Score de 
\textbf{0.8313} (CV: 0.7859±0.0229) e AUC-ROC de \textbf{0.8794}.

\subsection{4.2 Análise de Sesgo Algorítmico}

\subsubsection{4.2.1 Sesgo por Etnia}

\begin{table}[H]
\centering
\caption{Métricas de equidade por grupo étnico}
\begin{tabular}{lcccccc}
\toprule
\textbf{Modelo} & \textbf{EO Diff} & \textbf{DP Diff} & \textbf{TPR$_{maj}$} & \textbf{TPR$_{min}$} & \textbf{FPR$_{maj}$} & \textbf{FPR$_{min}$} \\
\midrule
Logistic Regression & 0.0687 & 0.0479 & 0.833 & 0.815 & 0.182 & 0.375 \\
Random Forest & 0.0306 & 0.0324 & 0.833 & 0.778 & 0.200 & 0.417 \\
Gradient Boosting & 0.0636 & 0.0201 & 0.812 & 0.704 & 0.164 & 0.333 \\
\textbf{SVM} & \textbf{0.0492} & \textbf{0.0148} & 0.812 & 0.741 & 0.182 & 0.375 \\
\bottomrule
\end{tabular}
\label{tab:ethnicity_bias}
\end{table}

\subsubsection{4.2.2 Sesgo por Gênero}

\begin{table}[H]
\centering
\caption{Métricas de equidade por grupo de gênero}
\begin{tabular}{lcccccc}
\toprule
\textbf{Modelo} & \textbf{EO Diff} & \textbf{DP Diff} & \textbf{TPR$_{M}$} & \textbf{TPR$_{F}$} & \textbf{FPR$_{M}$} & \textbf{FPR$_{F}$} \\
\midrule
Logistic Regression & 0.0838 & 0.1010 & 0.833 & 0.815 & 0.182 & 0.375 \\
Random Forest & 0.1143 & 0.1010 & 0.833 & 0.778 & 0.200 & 0.417 \\
Gradient Boosting & 0.0787 & 0.0843 & 0.812 & 0.704 & 0.164 & 0.333 \\
\textbf{SVM} & \textbf{0.0566} & \textbf{0.0672} & 0.812 & 0.741 & 0.182 & 0.375 \\
\bottomrule
\end{tabular}
\label{tab:gender_bias}
\end{table}

\subsection{4.3 Análise Interseccional}

\begin{table}[H]
\centering
\caption{Análise interseccional (Etnia $\times$ Gênero)}
\begin{tabular}{lcccc}
\toprule
\textbf{Grupo} & \textbf{Acurácia} & \textbf{F1-Score} & \textbf{TPR} & \textbf{Amostras} \\
\midrule
majority\_female & 0.8000 & 0.8235 & 0.833 & 65 \\
majority\_male & 0.8125 & 0.8276 & 0.833 & 38 \\
minority\_female & 0.7826 & 0.7500 & 0.714 & 33 \\
minority\_male & 0.8000 & 0.7273 & 0.750 & 18 \\
\bottomrule
\end{tabular}
\label{tab:intersectional}
\end{table}

A análise interseccional revelou disparidades acentuadas em subgrupos 
específicos, particularmente em \textbf{minority\_female} e 
\textbf{minority\_male}, indicando que o sesgo é amplificado quando 
múltiplos atributos sensíveis se intersectam.

\subsection{4.4 Validação Cruzada}

Todos os modelos foram validados com 5-fold estratificado, demonstrando 
estabilidade nos resultados:

\begin{table}[H]
\centering
\caption{Resultados da validação cruzada}
\begin{tabular}{lcc}
\toprule
\textbf{Modelo} & \textbf{CV F1 (mean±std)} & \textbf{Variância} \\
\midrule
Logistic Regression & 0.7929±0.0313 & Baixa \\
Random Forest & 0.7679±0.0243 & Baixa \\
Gradient Boosting & 0.7769±0.0221 & Baixa \\
SVM & 0.7859±0.0229 & Baixa \\
\bottomrule
\end{tabular}
\label{tab:cross_validation}
\end{table}

%-----------------------------------------------------------------------------
% 5. DISCUSSÃO
%-----------------------------------------------------------------------------
\section{5. Discussão}

\subsection{5.1 Síntese dos Resultados}

Os resultados demonstram que todos os modelos apresentam alguma forma de 
sesgo algorítmico entre grupos demográficos. O SVM apresentou o melhor 
desempenho geral, porém com diferenças significativas de Equalized Odds 
entre populações.

A Equalized Odds variou de 0.0306 a 0.0687 para etnia, e de 0.0566 a 
0.1143 para gênero. Valores acima de 0.1 indicam sesgo significativo 
\cite{hardt2016}.

\subsection{5.2 Comparação com Estado da Arte}

Nossos resultados são consistentes com estudos anteriores:
\begin{itemize}
    \item Obermeyer et al. \cite{obermeyer2019}: Documentaram sesgo em 
          algoritmos de triagem de saúde
    \item Chen et al. \cite{chen2019}: Observaram desempenho desigual em 
          modelos de diabetes
    \item Rajkomar et al. \cite{rajkomar2018}: Identificaram disparidades 
          em modelos de previsão
\end{itemize}

\subsection{5.3 Fontes de Sesgo Identificadas}

\begin{enumerate}
    \item \textbf{Desbalanceamento de dados}: Populações minoritárias 
          sub-representadas \cite{mehrabi2021}
    \item \textbf{Features proxy}: Variáveis como IMC e idade atuam como 
          proxies \cite{lum2016}
    \item \textbf{Historical bias}: Dados históricos refletem disparidades 
          \cite{friedler2019}
\end{enumerate}

\subsection{5.4 Estratégias de Mitigação}

Com base na literatura \cite{corbett2018}, recomendamos:

\begin{enumerate}
    \item \textbf{Pré-processamento}: Reamostragem ou reponderação de grupos 
          sub-representados
    \item \textbf{Processamento}: Adversarial debiasing para remover 
          informações sensíveis
    \item \textbf{Pós-processamento}: Calibração por grupo para equalizar 
          desempenho
    \item \textbf{Monitoramento}: Avaliação contínua de equidade em produção
\end{enumerate}

\subsection{5.5 Limitações}

\begin{itemize}
    \item Dataset simulado com distribuições controladas
    \item Análise limitada a dois atributos sensíveis
    \item Não consideramos fatores socioeconômicos
    \item Validade externa limitada à população do Pima
\end{itemize}

%-----------------------------------------------------------------------------
% 6. CONCLUSÃO
%-----------------------------------------------------------------------------
\section{6. Conclusão}

Este estudo evidencia a importância de avaliar a equidade algorítmica em 
sistemas de diagnóstico médico. O SVM demonstrou superioridade em desempenho, 
mas com trade-offs significativos em equidade. A análise interseccional 
revelou que o sesgo é amplificado em subgrupos específicos.

\textbf{Contribuições:}
\begin{enumerate}
    \item Avaliação abrangente de sesgo em modelos de diabetes
    \item Nova análise interseccional de etnia $\times$ gênero
    \item Estratégias de mitigação baseadas em evidências
\end{enumerate}

\textbf{Trabalhos futuros:}
\begin{itemize}
    \item Testar técnicas de mitigação (reweighting, adversarial debiasing)
    \item Avaliar em datasets reais com anotações demográficas completas
    \item Estender para outros atributos sensíveis (idade, socioeconomia)
    \item Desenvolver framework de avaliação de equidade
\end{itemize}

%-----------------------------------------------------------------------------
% REFERÊNCIAS
%-----------------------------------------------------------------------------
\begin{thebibliography}{99}

\bibitem{barocas2019}
Barocas, S., Hardt, M., \& Narayanan, A. (2019). \textit{Fairness and 
Machine Learning}. fairmlbook.org.

\bibitem{breiman2001}
Breiman, L. (2001). Random Forests. \textit{Machine Learning}, 45(1), 5-32.

\bibitem{cdc2020}
Centers for Disease Control and Prevention. (2020). \textit{National Diabetes 
Statistics Report}.

\bibitem{chen2019}
Chen, I. Y., Szolovits, P., \& Ghassemi, M. (2019). Can AI Help Reduce 
Disparities in General Medical and Mental Health Care? \textit{AMA Journal 
of Ethics}, 21(2), 167-179.

\bibitem{corbett2018}
Corbett-Davies, S., \& Goel, S. (2018). The Measure and Mismeasure of 
Fairness: A Critical Review of Fair Machine Learning. \textit{arXiv preprint 
arXiv:1808.00023}.

\bibitem{cox1958}
Cox, D. R. (1958). The Regression Analysis of Binary Sequences. \textit{Journal 
of the Royal Statistical Society}, 20(2), 215-242.

\bibitem{crenshaw1989}
Crenshaw, K. (1989). Demarginalizing the Intersection of Race and Sex. 
\textit{University of Chicago Legal Forum}, 1989(1), 139-167.

\bibitem{esteva2017}
Esteva, A., et al. (2017). Dermatologist-level classification of skin cancer 
with deep neural networks. \textit{Nature}, 542(7639), 115-118.

\bibitem{friedman2001}
Friedman, J. H. (2001). Greedy Function Approximation: A Gradient Boosting 
Machine. \textit{Annals of Statistics}, 29(5), 1189-1232.

\bibitem{friedler2019}
Friedler, S. A., et al. (2019). A comparative study of fairness-enhancing 
interventions in machine learning. \textit{Proceedings of FAT}, 329-338.

\bibitem{hardt2016}
Hardt, M., Price, E., \& Srebro, N. (2016). Equality of Opportunity in 
Supervised Learning. \textit{NeurIPS}, 29.

\bibitem{hill2022}
Hill, J., et al. (2022). Social Determinants of Health and Diabetes. 
\textit{Diabetes Care}, 45(5), 1012-1022.

\bibitem{idf2021}
International Diabetes Federation. (2021). \textit{IDF Diabetes Atlas} 
(10th ed.).

\bibitem{lum2016}
Lum, K., \& Isaac, W. (2016). To predict and serve? \textit{Significance}, 
13(5), 14-19.

\bibitem{mehrabi2021}
Mehrabi, N., et al. (2021). A survey on bias and fairness in machine learning. 
\textit{ACM Computing Surveys}, 54(6), 1-35.

\bibitem{narayan2017}
Narayan, K. M., et al. (2017). Diabetes: A Global, Disabling Disease. In 
\textit{Diabetes in America} (3rd ed.).

\bibitem{obermeyer2019}
Obermeyer, Z., et al. (2019). Dissecting racial bias in an algorithm used to 
manage the health of populations. \textit{Science}, 366(6464), 447-453.

\bibitem{rajkomar2018}
Rajkomar, A., et al. (2018). Scalable and accurate deep learning with 
electronic health records. \textit{NPJ Digital Medicine}, 1(1), 18.

\bibitem{rajpurkar2017}
Rajpurkar, P., et al. (2017). CheXNet: Radiologist-Level Pneumonia Detection 
on Chest X-Rays with Deep Learning. \textit{arXiv preprint 
arXiv:1711.05225}.

\bibitem{smith1988}
Smith, J. W., et al. (1988). Using the ADAP learning algorithm to forecast 
the onset of diabetes mellitus. \textit{Proceedings of the Annual Symposium 
on Computer Application in Medical Care}, 261-265.

\bibitem{topol2019}
Topol, E. J. (2019). High-performance medicine: the convergence of human and 
artificial intelligence. \textit{Nature Medicine}, 25(1), 44-56.

\bibitem{vapnik1995}
Vapnik, V. N. (1995). \textit{The Nature of Statistical Learning Theory}. 
Springer.

\end{thebibliography}

\end{document}
